%!TeX root=MemoriaTFG.tex

\chapter{Entorn d'Aprenentatge per Reforç}
\label{cap:entorns}

Per entrenar agents d'\ac{AR} al Truc cal adaptar el motor de joc a la interfície
estàndard que esperen les biblioteques d'\ac{AR}. Aquesta adaptació no és trivial:
el motor gestiona el torn de múltiples jugadors, manté un estat intern ric i emet
senyals de recompensa que les biblioteques genèriques no entenen de manera nativa.
Aquest capítol descriu com s'ha organitzat aquesta adaptació en capes, quines
decisions s'han pres per representar l'observació i les accions, i com s'ha dissenyat
el sistema de recompenses.

\section{Arquitectura en capes}
\label{sec:capes-entorn}

L'entorn s'organitza en dos nivells d'abstracció superposats (figura~\ref{fig:capes-entorn}):

\begin{enumerate}
  \item \textbf{Capa RLCard} (\texttt{TrucEnv} / \texttt{TrucEnvMa}): fa de pont
        entre el motor de joc i el món de l'\ac{AR}. Tradueix l'estat intern del motor
        en observacions estructurades, determina quines accions són legals a cada pas
        i recull les trajectòries completes de la partida.

  \item \textbf{Capa Gymnasium} (\texttt{TrucGymEnv} / \texttt{TrucGymEnvMa}): adapta
        la capa anterior a la interfície estàndard de Gymnasium (reinici i pas) perquè
        les biblioteques d'\acf{SB3} la puguin entrenar sense saber res del Truc. A
        més, fa transparent el torn del rival (avança el joc fins que torna a tocar a
        l'agent aprenent) i acumula les recompenses intermèdies.
\end{enumerate}

\begin{figure}[H]
  \centering
  \begin{tikzpicture}[
      box/.style={draw, rounded corners=3pt, minimum width=2.4cm, minimum height=0.9cm,
                  align=center, font=\small},
      arr/.style={-{Stealth}, thick},
      node distance=1.1cm and 0.8cm
  ]
    \node[box] (motor)   {\texttt{TrucGame}\\(motor)};
    \node[box, right=of motor] (rlcard)  {\texttt{TrucEnv}\\(RLCard)};
    \node[box, right=of rlcard] (gym)    {\texttt{TrucGymEnv}\\(Gymnasium)};
    \node[box, right=of gym]    (sb3)    {\ac{SB3} / PPO\\/ DQN / NFSP};

    \draw[arr] (motor)  -- (rlcard);
    \draw[arr] (rlcard) -- (gym);
    \draw[arr] (gym)    -- (sb3);

    \node[below=0.35cm of motor,  font=\footnotesize, text=gray] {lògica joc};
    \node[below=0.35cm of rlcard, font=\footnotesize, text=gray] {trajectòries};
    \node[below=0.35cm of gym,    font=\footnotesize, text=gray] {interfície RL};
    \node[below=0.35cm of sb3,    font=\footnotesize, text=gray] {algoritmes};
  \end{tikzpicture}
  \caption{Capes de l'entorn. \texttt{TrucEnvMa} i \texttt{TrucGymEnvMa} segueixen
           la mateixa estructura però amb episodi = 1~mà.}
  \label{fig:capes-entorn}
\end{figure}

La distinció entre \texttt{TrucEnv} i \texttt{TrucEnvMa} rau en la granularitat de
l'episodi: \texttt{TrucEnv} fa córrer una partida completa fins als 24~punts (centenars
de passos), mentre que \texttt{TrucEnvMa} atura l'episodi quan acaba la mà actual
(entre 6 i 20~passos aproximadament). Aquesta variant de gra fi és clau per al
\emph{curriculum learning}: l'agent aprèn primer a gestionar mans individuals i
posteriorment s'enfronta a partides senceres. El detall del curriculum es descriu
al capítol~\ref{cap:metodologia}.

A més d'aquests dos parells, hi ha dues variants especialitzades:

\begin{itemize}
  \item \textbf{Entorn de sessió} (\texttt{TrucGymEnvSessio}): l'episodi agrupa $N$
        mans consecutives contra el mateix oponent (per defecte $N=5$), que es
        selecciona del \emph{pool} a l'inici de cada episodi. En no reiniciar l'estat
        entre mans, obliga l'agent a acumular context al llarg de la sessió i li dona
        l'oportunitat de detectar el patró de joc del rival. S'usa a la Fase~4
        (memòria entre mans amb \texttt{RecurrentPPO}).

  \item \textbf{Entorn paral·lel}: executa diverses instàncies de l'entorn per mà en
        processos separats i en recull les trajectòries (amb les recompenses
        intermèdies) al procés principal. Permet col·lectar dades de molts episodis
        simultàniament, cosa que accelera l'entrenament propi de \ac{NFSP}.
\end{itemize}

L'entorn permet, a més, substituir el rival en calent sense reconstruir-lo. Aquesta
capacitat és imprescindible perquè els mecanismes d'entrenament puguin anar canviant
d'oponent durant una mateixa execució, tal com requereixen el \emph{pool} d'oponents
i el \emph{curriculum learning}.

\section{Espai d'observació}
\label{sec:observacio}

L'observació que rep l'agent en cada pas combina dues representacions
complementàries:

\begin{itemize}
  \item \texttt{obs\_cartes}: tensor de forma $(6, 4, 9)$, que codifica en
        format \emph{one-hot} la presència de cada carta en sis canals semàntics.
  \item \texttt{obs\_context}: vector de 24~valors reals que recull informació
        numèrica de l'estat de la partida.
\end{itemize}

Junts representen l'estat observable del joc en 240~dimensions. La funció
\texttt{flatten\_obs} és l'única font de veritat per a la conversió a vector pla:

\begin{equation}
  \mathbf{o} = \bigl[\,\texttt{obs\_cartes.flatten()}\;\|\;\texttt{obs\_context}\,\bigr]
  \in \mathbb{R}^{216+24} = \mathbb{R}^{240}
  \label{eq:obs-flat}
\end{equation}

Qualsevol codi que necessiti el vector pla (algoritmes SB3, NFSP, avaluació) crida
aquesta funció per garantir coherència.

\subsection*{Tensor de cartes (6, 4, 9)}

Les dimensions del tensor corresponen a: eix~0 = canal semàntic (6), eix~1 = pal
(4: Espases, Copes, Ors, Bastos), eix~2 = rang (9: $\{1,3,4,5,6,7,10,11,12\}$).
Cada cel·la val~1 si la carta corresponent hi és present i~0 altrament.

\begin{table}[H]
  \centering
  \begin{tabular}{cll}
    \toprule
    Canal & Contingut & Descripció \\
    \midrule
    0 & Mà actual         & Cartes que té l'agent a la mà \\
    1 & Historial propi   & Cartes que ja ha jugat l'agent \\
    2 & Historial Rival 1 & Cartes jugades pel rival de la dreta \\
    3 & Historial Company & Cartes jugades pel company (2v2) \\
    4 & Historial Rival 2 & Cartes jugades pel segon rival (2v2) \\
    5 & Senyes company    & Cartes revelades per les senyes del company \\
    \bottomrule
  \end{tabular}
  \caption{Els sis canals del tensor de cartes \texttt{obs\_cartes}.}
  \label{tab:canals-cartes}
\end{table}

Els índexs de canal 1–4 s'assignen de manera relativa al jugador actual: l'offset
$k$ correspon al jugador $(id + k) \bmod n$, de manera que la mateixa estructura de
tensor és vàlida per a qualsevol posició de la taula i per a les modalitats 2v2 i 1v1.

\subsection*{Vector de context (24 dimensions)}

\begin{table}[H]
  \centering
  \begin{tabular}{clll}
    \toprule
    Índex(os) & Significat & Rang & Codificació \\
    \midrule
    0   & Puntuació equip propi    & $[0,1]$ & $p_{\text{propi}} / 24$ \\
    1   & Puntuació equip rival    & $[0,1]$ & $p_{\text{rival}} / 24$ \\
    2   & Nivell aposta de Truc    & $[0,1]$ & $\text{level}_{\text{truc}} / 24$ \\
    3   & Nivell aposta d'Envit    & $[0,1]$ & $\text{level}_{\text{envit}} / 24$ \\
    4   & Fase del torn            & $\{0,1\}$ & 0 = senyes, 1 = apostes/joc \\
    5   & Ronda actual             & $[0,1]$ & $\text{ronda} / \text{cartes\_jugador}$ \\
    6–9 & Qui és \emph{mà}         & $\{0,1\}^4$ & One-hot relatiu a l'agent \\
    10–13 & Qui ha cantat Truc     & $\{0,1\}^4$ & One-hot relatiu; 0 si no s'ha cantat \\
    14–16 & Qui ha cantat Envit    & $\{0,1\}^3$ & One-hot relatiu; 0 si no s'ha cantat \\
    17  & Rondes guanyades (propi) & $[0,1]$ & $r_{\text{propi}} / 3$ \\
    18  & Rondes guanyades (rival) & $[0,1]$ & $r_{\text{rival}} / 3$ \\
    19  & Guanyador ronda 1        & $\{-1,0,1\}$ & $+1$ = propi, $-1$ = rival, $0$ = empat \\
    20  & Guanyador ronda 2        & $\{-1,0,1\}$ & $+1$ = propi, $-1$ = rival, $0$ = empat \\
    21  & Envit acceptat           & $\{0,1\}$ & Binari \\
    22  & Truc pendent             & $\{0,1\}$ & 1 si \texttt{TRUC\_PENDING} \\
    23  & Envit pendent            & $\{0,1\}$ & 1 si \texttt{ENVIT\_PENDING} \\
    \bottomrule
  \end{tabular}
  \caption{Les 24 dimensions del vector de context \texttt{obs\_context}.}
  \label{tab:context-obs}
\end{table}

Totes les variables s'expressen de manera relativa a la perspectiva de l'agent
actual. Això fa que el tensor sigui invariant a la posició de la taula: no cal
cap transformació addicional per canviar de jugador.

Aquest vector de 240~dimensions és l'entrada de la xarxa de l'agent. L'arquitectura
concreta de l'extractor de característiques que el processa (en particular la
comparació entre un \acf{MLP} pla i un extractor convolucional que respecta
l'estructura $6 \times 4 \times 9$ del tensor de cartes) es tracta al
capítol~\ref{cap:metodologia}.

\section{Espai d'accions i màscara}
\label{sec:accions-entorn}

L'espai d'accions conté 19 accions discretes, definides al capítol~\ref{cap:joc}. No
totes elles són legals en tot moment: si hi ha un Truc pendent, l'agent només pot
respondre (vull, fora o contra); si s'estan jugant cartes no pot apostar si ja ha
apostat abans en la mateixa mà, etc.

A cada pas, l'entorn calcula el conjunt d'accions vàlides segons la situació i el posa
a disposició de l'algorisme de dues maneres complementàries:

\begin{itemize}
  \item \textbf{Correcció d'accions il·legals}: si el model tria una acció que no és
        vàlida (cosa habitual a les primeres fases, quan la política és gairebé
        aleatòria), l'entorn la reemplaça per una acció legal en lloc de generar un
        error. Així l'entrenament mai no s'interromp per una jugada invàlida.
  \item \textbf{Màscara explícita}: amb els algorismes que ho admeten (com el \ac{PPO}
        amb màscara d'accions de \ac{SB3}), el conjunt legal es passa directament a la
        política perquè assigni probabilitat zero a les accions il·legals.
\end{itemize}

En jocs amb apostes encadenades com el Truc, la màscara és especialment important:
sense ella, l'agent podria intentar re-apostar quan ja no és possible, o respondre
a una aposta que ja ha estat resolta. La separació entre les accions definides al
motor i les legals en cada pas permet que el motor mantingui la lògica de negoci
sense barrejar-la amb la política de l'agent.

\section{Reward shaping}
\label{sec:reward}

En una partida fins a 24~punts, la recompensa final pot trigar centenars de passos
a arribar. Sense senyals intermèdies, l'agent sabria si ha guanyat o perdut però no
podria atribuir el resultat a cap acció concreta: el problema d'assignació de crèdit
fa que aprendre a gestionar l'envit o a forçar el truc sigui pràcticament impossible
en condicions de recompensa escassa. La solució és el \emph{reward shaping}.

\subsection*{Recompensa terminal}

Quan acaba la partida (o la mà, en la variant de gra fi) l'entorn lliura la
recompensa terminal:

\begin{itemize}
  \item \textbf{\texttt{TrucEnv} (partida completa)}: $+1{,}0$ per a l'equip
        guanyador, $-1{,}0$ per al perdedor.
  \item \textbf{\texttt{TrucEnvMa} (mà individual)}: $\pm\,\text{punts}/24$,
        on \emph{punts} són els punts de la mà (Truc + Envit). Normalitzar per~24
        manté la recompensa a l'escala de la partida completa.
\end{itemize}

\subsection*{Recompenses intermèdies}

El motor associa a cada transició una recompensa intermèdia que reflecteix el
resultat immediat de l'acció. El tractament difereix segons la capa. A la capa
RLCard, que per disseny deixa totes les transicions intermèdies a zero, un postprocés
recorre la trajectòria un cop acabat l'episodi i substitueix aquests zeros pel valor
acumulat de l'equip, deixant intacte el \emph{payoff} final. A la capa Gymnasium, en
canvi, la recompensa intermèdia s'extreu i s'aplica directament a cada pas,
amplificada per un factor~5:

\begin{equation}
  r_t = \texttt{reward\_intermedis}[\text{equip}] \times 5{,}0
  \label{eq:reward-gym}
\end{equation}

El factor~5 equilibra l'escala entre els senyals intermedis (ordre de~0.1--0.5) i
la recompensa terminal (±1.0), donant suficient pes al feedback immediat sense
eclipsar l'objectiu final.

Les recompenses intermèdies que emet el motor es desglossen de la manera següent:

\begin{itemize}
  \item \textbf{Guanyar una ronda}: $+0{,}40 \times w_r$ per al guanyador i
        $-0{,}20 \times w_r$ per al perdedor, on $w_r$ és un pes decreixent
        per ronda ($w_1=2{,}0$, $w_2=1{,}0$, $w_3=0{,}5$) que reflecteix que
        les primeres rondes són més determinants. Si perdre la ronda implica
        perdre la mà, la penalització s'amplifica per un factor~$1{,}5$.

  \item \textbf{Envit acceptat}: recompensa proporcional als punts en joc per
        a l'equip que matemàticament guanya l'envit.

  \item \textbf{Envit rebutjat (\texttt{fora\_envit})}: $+1{,}5\times$ per a
        l'equip que força el rebuig, $-1{,}0\times$ per al que es retira.

  \item \textbf{Truc rebutjat (\texttt{fora\_truc} en resposta)}: $+1{,}5\times$
        per a l'agressiu, $-1{,}0\times$ per al retirat.

  \item \textbf{Retirada voluntària (\texttt{fora\_truc} sense resposta pendent)}:
        $+1{,}0\times$ per a l'oponent, $-1{,}2\times$ per al que fuig, per
        penalitzar les fugides injustificades.

  \item \textbf{Guanyar la mà (truc net)}: $\pm 1{,}0\times$ el factor base.
\end{itemize}

La variant per mans simplifica el sistema: acumula un únic reward net al final de
cada mà com $\pm\,\text{punts}/24$, sense senyals de ronda individuals. Això afavoreix
un aprenentatge més estable per a la fase de \emph{curriculum learning} on cada
episodi és una única mà.
